% GENERATED FILE. Edit canonical lessons, then run npm run generate:books.
% canonical-source-hash: fnv1a64:fd95be8b2f93e9b8
% canonical-lessons: GE-C260-wenn-falls, GE-C260-wenn-es-regnet, GE-C260-es-kommt-darauf-an, GE-C260-review

\chapter{If — wenn, falls, es kommt darauf an}
\label{ch:ge-if}
\hlchaptermodality{260}

\begin{chapteropening}
\textbf{By the end of this chapter:} \emph{I can state a real condition with wenn or falls and the present, and say that something depends, naming on what.}

Complete GE-C260-review: a plan with a branch in it, a real condition with wenn or falls, and an answer that depends.
\end{chapteropening}

\section[wenn, falls]{wenn, falls — if}

\label{lesson:GE-C260-wenn-falls}



\begin{quote}
Every time I was in Berlin, I visited my grandmother. Say it with \emph{wenn}. (\emph{Wenn ich in Berlin war, besuchte ich meine Großmutter.})

That \emph{wenn} was \textquotedblleft{}whenever\textquotedblright{}. The same word also means \textquotedblleft{}if\textquotedblright{}.
\end{quote}

\begin{grammarlens}[title={Grammar Lens: wenn for if}]
German's everyday \textquotedblleft{}if\textquotedblright{} is \textbf{wenn}, and its verb goes to the end as before:

\begin{quote}
\emph{Wenn du Zeit} \textbf{hast}, \emph{komm!} — If you have time, come!
\end{quote}

English speakers hear \textquotedblleft{}when\textquotedblright{} in it at first, and often that is fine: \textquotedblleft{}if\textquotedblright{} and \textquotedblleft{}whenever\textquotedblright{} sit close together.

When you need it to be only \textquotedblleft{}if\textquotedblright{}, \textquotedblleft{}in case it happens\textquotedblright{}, German has \textbf{falls}:

\begin{quote}
\emph{Falls der Zug Verspätung} \textbf{hat}, \emph{rufe ich dich an.} — If the train is late, I'll call you.
\end{quote}

\emph{Falls} does not mean \textquotedblleft{}when\textquotedblright{}. It is a little more careful than \emph{wenn}, and it suits a plan with a fallback in it. Both send the verb to the end.
\end{grammarlens}

\begin{cousinweb}
\textbf{falls} is built on \textbf{der Fall}, \textquotedblleft{}the case\textquotedblright{}, which EtymDB 2.1 traces through Middle High German \emph{val} to Old High German \emph{val}, \textquotedblleft{}a fall\textquotedblright{}; it belongs with the verb \textbf{fallen}. The \emph{-s} is the old adverb ending of \emph{rechts} and \emph{links}, so \emph{falls} is \textquotedblleft{}in case\textquotedblright{}, by the usual account grown out of phrases like \emph{im Falle}, \textquotedblleft{}in the case\textquotedblright{}.

\textbf{wenn} is the old twin of \emph{wann}, from chapter two hundred and fifty-two.
\end{cousinweb}

\subsection*{Your turn}
Say these aloud:

\begin{itemize}
\raggedright
  \item \emph{Wenn du Zeit hast, komm!}
  \item \emph{Falls der Zug Verspätung hat, rufe ich dich an.}
  \item the verb last in both: \emph{hast}, \emph{hat}
  \item \emph{der Fall}, the case; \emph{falls}, in case
  \item one of your own with \emph{falls}, about tomorrow
\end{itemize}

\begin{tcolorbox}[breakable,colback=teal!4,colframe=teal!35!black,title={Before you move on}]
\textquotedblleft{}If you have time\textquotedblright{}? (\emph{\textbf{Wenn du Zeit hast.}}) Which word means only \textquotedblleft{}if, in case\textquotedblright{}, and not \textquotedblleft{}when\textquotedblright{}? (\emph{\textbf{Falls.}}) Which noun is inside it? (\emph{\textbf{Der Fall}}, the case.)
\end{tcolorbox}
\section[wenn es regnet, …]{wenn es regnet — a real condition}

\label{lesson:GE-C260-wenn-es-regnet}



\begin{quote}
Say \textquotedblleft{}it's raining\textquotedblright{}. (\emph{Es regnet.}) Then \textquotedblleft{}if you have time\textquotedblright{}. (\emph{Wenn du Zeit hast.})
\end{quote}

\begin{grammarlens}[title={Grammar Lens: the present, for tomorrow too}]
A real condition is one that may well come true. German puts both halves in the present, even when they are about tomorrow:

\begin{quote}
\emph{Wenn es morgen regnet, fahren wir nicht.} — If it rains tomorrow, we won't go. \emph{Wenn ich müde bin, nehme ich ein Taxi.} — If I'm tired, I take a taxi.
\end{quote}

English does the same in the \emph{if} half: \textquotedblleft{}if it rains\textquotedblright{}. German usually goes one step further and keeps the other half in the present as well, where English would say \textquotedblleft{}we won't go\textquotedblright{}. The word \emph{morgen} already tells you it is tomorrow.
\end{grammarlens}

\begin{culture}
You heard it first with \emph{als}. When the \emph{wenn} clause comes first, it fills the first place of the whole sentence, so the main verb comes next, right after the comma:

\begin{quote}
\emph{Wenn es morgen} \textbf{regnet}, \textbf{fahren} \emph{wir nicht.}
\end{quote}

Two verbs side by side, one ending the condition, one opening the result.

Turn it round, and each clause keeps its own order:

\begin{quote}
\emph{Wir} \textbf{fahren} \emph{nicht, wenn es morgen} \textbf{regnet}.
\end{quote}

Some speakers add \emph{dann} to the result. It takes the place after the comma, and the verb still comes second: \emph{Wenn es regnet, dann fahren wir nicht.}
\end{culture}

\begin{cousinweb}
\textbf{regnen} is built on \textbf{der Regen}, from Old High German \emph{regan}, which EtymDB 2.1 sets beside English \emph{rain}: one inherited word. \emph{Es regnet} and \textquotedblleft{}it rains\textquotedblright{} both need an empty \emph{es}, \emph{it}, because nobody does the raining.
\end{cousinweb}

\subsection*{Your turn}
Say these aloud:

\begin{itemize}
\raggedright
  \item \emph{Wenn es morgen regnet, fahren wir nicht.}
  \item the two verbs at the comma: \emph{regnet, fahren}
  \item turned round: \emph{Wir fahren nicht, wenn es morgen regnet.}
  \item \emph{Wenn ich müde bin, nehme ich ein Taxi.}
  \item with \emph{falls}: \emph{Falls es regnet, bleiben wir in der Wohnung.}
\end{itemize}

\begin{tcolorbox}[breakable,colback=teal!4,colframe=teal!35!black,title={Before you move on}]
\textquotedblleft{}If it rains tomorrow, we won't go\textquotedblright{}? (\emph{\textbf{Wenn es morgen regnet, fahren wir nicht.}}) Which tense for tomorrow, in both halves? (\textbf{The present.}) After a \emph{wenn} clause, what comes straight after the comma? (\textbf{The main verb.})
\end{tcolorbox}
\section[es kommt darauf an / das hängt davon ab]{es kommt darauf an — it depends}

\label{lesson:GE-C260-es-kommt-darauf-an}



\begin{quote}
Say \textquotedblleft{}if it rains tomorrow, we won't go\textquotedblright{}. (\emph{Wenn es morgen regnet, fahren wir nicht.})

Now someone asks if you are coming, and the honest answer is \textquotedblleft{}it depends\textquotedblright{}.
\end{quote}

\subsection*{You'll want to know: es kommt darauf an}
\textbf{Es kommt darauf an.} — \textquotedblleft{}It depends\textquotedblright{}: \emph{ess kommt dah-ROWF ahn}.

\begin{quote}
\emph{Kommst du morgen?} — \emph{Es kommt darauf an.}
\end{quote}

It is \emph{ankommen}, split as usual, with \emph{darauf}, \textquotedblleft{}on that\textquotedblright{}, in the middle. In quick speech it shrinks to \emph{kommt drauf an}.

To name the thing it depends on, swap \emph{darauf} for \emph{auf} and the thing:

\begin{quote}
\emph{Es kommt auf das Wetter an.} — It depends on the weather.
\end{quote}

And to ask \textquotedblleft{}on what?\textquotedblright{}, German builds the question the way it built \emph{warum}: \textbf{worauf?}

\begin{grammarlens}[title={Grammar Lens: das hängt davon ab, ob …}]
The second way is \textbf{das hängt davon ab}, \textquotedblleft{}that depends on it\textquotedblright{}. It is \emph{abhängen}, to hang off, split like \emph{ankommen}:

\begin{quote}
\emph{Das hängt vom Wetter ab.} — That depends on the weather.
\end{quote}

With \emph{abhängen} the thing comes after \textbf{von}: \emph{vom Wetter}, from \emph{von dem}.

When what it depends on is a whole question, German opens it with \textbf{ob}, \textquotedblleft{}whether\textquotedblright{}, and the verb goes to the end:

\begin{quote}
\emph{Das hängt davon ab, ob es} \textbf{regnet}. — That depends on whether it rains. \emph{Es kommt darauf an, ob ich Zeit} \textbf{habe}. — It depends on whether I have time.
\end{quote}

\emph{Ob} is the \textquotedblleft{}if\textquotedblright{} of \textquotedblleft{}I don't know if\textquotedblright{}: it asks, where \emph{wenn} sets a condition.
\end{grammarlens}

\begin{cousinweb}
\textbf{ankommen} is \emph{an-} + \emph{kommen}, and EtymDB 2.1 sets English \emph{oncome}, from Old English \emph{oncuman}, beside it. \textbf{hängen} $\leftarrow$ Middle High German \emph{hengen} $\leftarrow$ Old High German \emph{hengen} (EtymDB 2.1). \emph{Abhängen}, to hang off, gives the same picture as English \emph{depend}, from Latin \emph{dēpendēre}, \textquotedblleft{}to hang down from\textquotedblright{}: what depends on a thing hangs from it.
\end{cousinweb}

\subsection*{Your turn}
Say these aloud:

\begin{itemize}
\raggedright
  \item \emph{Es kommt darauf an.}
  \item \emph{Es kommt auf das Wetter an.}
  \item \emph{Das hängt vom Wetter ab.}
  \item \emph{Das hängt davon ab, ob es regnet.}
  \item \emph{Kommst du morgen? Es kommt darauf an, ob ich Zeit habe.}
\end{itemize}

\begin{tcolorbox}[breakable,colback=teal!4,colframe=teal!35!black,title={Before you move on}]
\textquotedblleft{}It depends\textquotedblright{}? (\emph{\textbf{Es kommt darauf an.}}) \textquotedblleft{}That depends on the weather\textquotedblright{}? (\emph{\textbf{Das hängt vom Wetter ab.}}) Which word means \textquotedblleft{}whether\textquotedblright{}, and where does its verb go? (\emph{\textbf{Ob}}, \textbf{and the verb goes to the end.})
\end{tcolorbox}
\section[(review)]{Review — if, and it depends}

\label{lesson:GE-C260-review}



\begin{quote}
Say \textquotedblleft{}in case the train is late\textquotedblright{}. (\emph{Falls der Zug Verspätung hat.})
\end{quote}

\begin{grammarlens}[title={What you've built}]
Three lessons, and a plan can now branch. A real condition with \emph{wenn}, or with \emph{falls} when you mean only \textquotedblleft{}in case\textquotedblright{}. Both halves in the present, even for tomorrow. Verb, comma, verb when the condition comes first. And when you cannot promise, \emph{es kommt darauf an} or \emph{das hängt davon ab}, with \emph{ob} for the open question.

One exchange that uses most of it:

\begin{quote}
\emph{Kommst du morgen?} \emph{Es kommt darauf an, ob ich Zeit habe. Wenn ich Zeit habe, komme ich. Falls nicht, rufe ich dich an.}
\end{quote}

\emph{Falls nicht}, \textquotedblleft{}if not\textquotedblright{}, is a whole condition in two words.
\end{grammarlens}

\subsection*{Your turn}
Say these aloud:

\begin{itemize}
\raggedright
  \item \emph{Wenn es morgen regnet, fahren wir nicht.}
  \item \emph{Falls der Zug Verspätung hat, rufe ich dich an.}
  \item \emph{Es kommt darauf an, ob ich Zeit habe.}
  \item \emph{Wenn ich Zeit habe, komme ich. Falls nicht, rufe ich dich an.}
  \item your own plan for the weekend, with one \emph{wenn} and one \emph{falls nicht}
  \item \emph{Es regnet seit zwei Tagen. Falls es morgen regnet, sollte man einen Regenschirm mitnehmen.}
\end{itemize}

\begin{tcolorbox}[breakable,colback=teal!4,colframe=teal!35!black,title={Before you move on}]
Which tense in both halves of a real condition about tomorrow? (\textbf{The present.}) \textquotedblleft{}If not\textquotedblright{}? (\emph{\textbf{Falls nicht.}}) \textquotedblleft{}That depends on whether it rains\textquotedblright{}? (\emph{\textbf{Das hängt davon ab, ob es regnet.}})
\end{tcolorbox}
